\section*{Capítulo \ref{ch:b2:weather-systems} --- Sistemas meteorológicos}

\begin{solution}{exo:b2:weather-systems:1}
(a) $f = 2 \times 7.292 \times 10^{-5} \times \sin\ang{50} = \qty{1.117e-4}{s^{-1}}$. (b) Com
$\bar T = \qty{278}{K}$, $\tan\alpha = 1.117 \times 10^{-4} \times 278 \times 15/(9.81 \times 6)
= 0.00791$, cerca de 1 para 126. (c) $1/0.00791 = 126$: uns \qty{126}{km} atrás da frente. (d) A
inclinação é inversamente proporcional ao contraste de temperatura: cai à metade, para 1 para 253.
\end{solution}

\begin{solution}{exo:b2:weather-systems:2}
(a) $6 \times 150 = \qty{900}{km}$ à frente. (b) $900/25 = \qty{36}{h}$, um dia e meio. (c) Primeiro
os cirros altos, depois os cirrostratos, depois os altostratos; a nuvem espessa-se e baixa, e a chuva
começa a cair da nuvem baixa e espessa à medida que a frente se aproxima.
\end{solution}

\begin{solution}{exo:b2:weather-systems:3}
Os símbolos apontam no sentido em que a frente se move, para o ar mais quente à frente de uma frente fria ou para o
ar mais frio à frente de uma frente quente. (a) Uma frente fria: o ar frio fica a noroeste dela, atrás
dela; ela se move para sudeste. (b) Uma frente quente: o ar frio fica ao norte dela, à sua frente; ela se move para
o norte. (c) Uma \omterm{def:b2:weather-systems:fronts}{frente oclusa}, que liga a baixa ao ponto triplo de onde partem as frentes fria e
quente; há ar frio dos dois lados dela, com o ar quente levantado acima; ela se move no sentido para o qual seus
símbolos apontam. (d) Uma \omterm{def:b2:weather-systems:fronts}{frente estacionária}: os triângulos apontam para o sul, para o ar quente, logo o ar
frio fica ao norte; ela quase não se move.
\end{solution}

\begin{solution}{exo:b2:weather-systems:4}
(a) $\beta = 2 \times 7.292 \times 10^{-5} \times \cos\ang{45}/(6.371 \times 10^6) =
\qty{1.62e-11}{m^{-1}.s^{-1}}$. (b) Para \qty{4000}{km}, $k = \qty{1.571e-6}{m^{-1}}$, $\beta/k^2 =
\qty{6.6}{m/s}$ e $c = 12 - 6.6 = \qty{5.4}{m/s}$; para \qty{8000}{km}, $\beta/k^2 =
\qty{26.3}{m/s}$ e $c = \qty{-14.3}{m/s}$. (c) A onda curta move-se para leste, \qty{470}{km} por dia;
a longa move-se para oeste, \qty{1240}{km} por dia, contra o vento.
\end{solution}

\begin{solution}{exo:b2:weather-systems:5}
(a) $c = 0$ quando $U = \beta/k^2$, logo $k = \sqrt{\beta/U}$ e $\lambda = 2\pi/k =
2\pi\sqrt{U/\beta}$. (b) $2\pi\sqrt{12/(1.62 \times 10^{-11})} = \qty{5410}{km}$. (c) O círculo
de latitude \ang{45} mede $2\pi \times 6371 \times \cos\ang{45} = \qty{28300}{km}$: cerca de 5.2
ondas, ou seja, cinco na prática. (d) $\lambda$ cresce como $\sqrt U$: um vento de oeste mais forte torna as
ondas estacionárias mais longas, e cabem menos delas em volta do hemisfério.
\end{solution}

\begin{solution}{exo:b2:weather-systems:6}
(a) Em $C_+ = (\sqrt{72}, \sqrt{72}, 27)$: $\dot x = 10(\sqrt{72} - \sqrt{72}) = 0$, $\dot y =
\sqrt{72}(28 - 27) - \sqrt{72} = 0$, $\dot z = 72 - \tfrac83 \times 27 = 72 - 72 = 0$.
(b) $\mathrm{e}^{-13.67} = 1.16 \times 10^{-6}$: o volume encolhe quase um milhão de vezes numa
unidade de tempo. (c) Um volume pode encolher enquanto é esticado numa direção e comprimido com muito mais
força nas outras: a distância ao longo da direção esticada cresce, e duas trajetórias vizinhas
separam-se ao longo dela, enquanto o volume entre elas se anula.
\end{solution}

\begin{solution}{exo:b2:weather-systems:7}
(a) $\ln 2/0.906 = 0.765$ unidades de tempo. (b) $\ln(10/10^{-8})/0.906 = \ln(10^9)/0.906 = 22.9$ unidades
de tempo. (c) $\ln 1000/0.906 = 7.6$ unidades de tempo: uma melhoria de mil vezes do estado inicial
compra apenas um terço do tempo encontrado em (b), pois $\ln 10^3/\ln 10^9 = 1/3$.
\end{solution}

\begin{solution}{exo:b2:weather-systems:8}
(a) $(303 - 195)/195 = 0.554$ e $V_{\max} = \sqrt{0.9 \times 0.554 \times 15\,000} =
\qty{86.5}{m/s}$. (b) Bem acima dos \qty{70}{m/s} da categoria 5, perto dos \qty{86}{m/s} que
cinco tempestades recentes ultrapassaram. (c) As tempestades raramente ficam muito tempo no estado estacionário que a teoria
supõe; o cisalhamento vertical do vento, o ar seco aspirado para a tempestade e o ar que desce à sua volta, que
a teoria deixa de fora, enfraquecem-nas.
\end{solution}

\begin{solution}{exo:b2:weather-systems:9}
(a) $(303 - 195)/303 = 0.356$. (b) $(303 - 195)/195 = 0.554$; a razão é $0.554/0.356 =
303/195 = 1.554$. (c) Numa máquina de Carnot o trabalho sai da máquina. Num ciclone tropical o
trabalho é dissipado por atrito perto da superfície do mar, e esse calor aquece o ar que entra à
temperatura $T_s$, onde a máquina recebe seu calor: ele é reutilizado, e a máquina comporta-se
como se fosse mais eficiente pelo fator $T_s/T_o$.
\end{solution}

\begin{solution}{exo:b2:weather-systems:10}
(a) A dedução da \cref{prop:b2:weather-systems:slope} sem aproximação dá
$\mathrm{d}z/\mathrm{d}x = f(\rho_c v_c - \rho_w v_w)/[g(\rho_c - \rho_w)]$. Um fluido mais denso
apoiado sobre um mais leve é instável: a menor perturbação o faz tombar. Numa frente em
repouso o ar frio, mais denso, tem portanto de ficar por baixo, que é a cunha da definição; então
$\rho_c > \rho_w$, a superfície sobe em direção ao ar frio, e $\rho_c v_c > \rho_w v_w$.
(b) $\zeta = \partial v/\partial x = 10/(50 \times 10^3) = \qty{2e-4}{s^{-1}}$, cerca de $1.9f$ com
$f = \qty{1.03e-4}{s^{-1}}$. (c) Positiva no Hemisfério Norte: ciclônica, e tão grande quanto
o próprio parâmetro de Coriolis; as frentes são faixas de forte vorticidade ciclônica.
\end{solution}

\begin{solution}{exo:b2:weather-systems:11}
(a) $\log_2(6/0.5) = \log_2 12 = 3.58$ duplicações, $3.58 \times 2.4 = 8.6$~dias. (b) Um
dia a mais de previsão exige um tempo de duplicação a menos de crescimento antes de o erro ficar inútil: o
erro inicial deve diminuir $2^{1/2.4} = 1.33$ vezes a cada década. (c) Uma redução por dez equivale a
$\log_2 10 = 3.32$ duplicações; se ela ganha apenas 4 a 5 dias, os pequenos erros duplicam em $4/3.32 = 1.2$
a $5/3.32 = 1.5$ dias: os erros menores crescem mais depressa que os maiores medidos nas
previsões dos anos 1980, através do calor latente liberado na convecção em pequenas escalas.
\end{solution}

\begin{solution}{exo:b2:weather-systems:12}
(a) A \qty{30}{\celsius}, $e_s = \qty{43.5}{hPa}$ e $q_s^* = \qty{27.1}{g/kg}$; o ar a
\qty{29}{\celsius} tem $e = 0.8 \times 41.0 = \qty{32.8}{hPa}$ e $q_a = \qty{20.4}{g/kg}$;
$h_o^* - h^* = 1004.5 \times 1 + 2.501 \times 10^6 \times 0.0067 = \qty{17.8}{kJ/kg}$, contra
\qty{16.8}{kJ/kg}. (b) $(303.15 - 200)/200 = 0.516$ e $V_{\max} = \sqrt{0.9 \times 0.516
\times 17\,800} = \qty{90.9}{m/s}$, contra \qty{87.9}{m/s}. (c) Cerca de \qty{3.4}{\percent} por
kelvin. O desequilíbrio aumenta \qty{6}{\percent}, o que acrescenta \qty{3}{\percent} à velocidade;
o fator $(T_s - T_o)/T_o$ aumenta \qty{1}{\percent}, o que acrescenta \qty{0.5}{\percent}: o
combustível, a umidade que o mar mais quente pode fornecer, é o que mais contribui, através do aumento exponencial de
$e_s$.
\end{solution}

\begin{solution}{pb:b2:weather-systems:1}
\textbf{1.} $(480 - 390)/6 = \qty{15}{km/h}$.
\textbf{2.} $390/15 = \qty{26}{h}$ depois das 12 UTC: às 14 UTC do dia seguinte.
\textbf{3.} $f = \qty{1.117e-4}{s^{-1}}$.
\textbf{4.} Com $\bar T = \qty{281}{K}$, $\tan\alpha = 1.117 \times 10^{-4} \times 281 \times
16/(9.81 \times 8) = 0.0064$, cerca de 1 para 156.
\textbf{5.} A frente está então $6 \times 15 = \qty{90}{km}$ além da cidade; a superfície está
$90\,000 \times 0.0064 = \qty{576}{m}$ acima dela.
\textbf{6.} $1500/0.0064 = \qty{234}{km}$ atrás da frente, atingidos $234/15 = \qty{15.6}{h}$ depois
da passagem.
\textbf{7.} Cúmulos e cumulonimbos, aguaceiros, talvez trovoada, rajadas e uma queda súbita da
temperatura; a pressão, que cai à medida que a frente se aproxima, começa a subir quando ela passa; depois
um ar mais fresco e mais limpo.
\textbf{8.} $\beta = 2 \times 7.292 \times 10^{-5} \times \cos\ang{50}/(6.371 \times 10^6) =
\qty{1.47e-11}{m^{-1}.s^{-1}}$.
\textbf{9.} $k = 2\pi/(5 \times 10^6) = \qty{1.257e-6}{m^{-1}}$ e $\beta/k^2 = \qty{9.3}{m/s}$.
\textbf{10.} $c = 20 - 9.3 = \qty{10.7}{m/s}$, para leste.
\textbf{11.} $10.7 \times 86\,400 = \qty{924}{km}$ para leste.
\textbf{12.} $\lambda_s = 2\pi\sqrt{20/(1.47 \times 10^{-11})} = \qty{7330}{km}$; o círculo de
latitude \ang{50} mede $2\pi \times 6371 \times \cos\ang{50} = \qty{25700}{km}$, cerca de 3.5
dessas ondas.
\textbf{13.} À frente dele, do seu lado leste, onde as ondas curtas induzem ar ascendente e a
pressão à superfície cai.
\textbf{14.} $\lambda = \ln(3.9 \times 10^{-3}/5 \times 10^{-7})/(14 - 4) = \ln 7800/10 = 0.896$.
\textbf{15.} $\ln 2/0.896 = 0.77$ unidades de tempo.
\textbf{16.} \qty{1}{\percent} abaixo de 0.906: a média de 200 pares, um tempo finito e leituras num
gráfico.
\textbf{17.} Quando as duas rodadas estão muito afastadas, ficam em pontos sem relação do atrator, e
sua distância não pode exceder o tamanho dele: o crescimento satura.
\textbf{18.} $\ln 10/0.896 = 2.6$ unidades de tempo mais tarde.
\textbf{19.} $\log_2(3/0.1) = \log_2 30 = 4.91$ duplicações.
\textbf{20.} $4.91 \times 2.4 = 11.8$~dias.
\textbf{21.} Um tempo de duplicação: 2.4~dias.
\textbf{22.} $4.91 \times 5 = 24.5$~dias: a estimativa antiga era otimista demais.
\textbf{23.} Com um tempo de duplicação de 2.4 dias, um erro dez vezes menor ganharia $3.32 \times 2.4
= 8.0$~dias, e não 4 a 5: os pequenos erros crescem mais depressa que 2.4 dias por duplicação, e uma
estimativa baseada nesse tempo de duplicação superestima o horizonte.
\textbf{24.} Esta previsão poderia ser útil por no máximo cerca de \textbf{12 dias}, e na prática
por menos.
\end{solution}
